\chapter{Interacción con la base de datos} \label{Anexo:InteraccionBBDD}
%\section{Cuestionario de usabilidad}

%%%%%%%%%%%%%%%%%%%%%%%%%%%%%%%% INTERACCIONES %%%%%%%%%%%%%%%%%%%%%%%%%%%%%%%%        
    Como se ha mencionado previamente, en el proyecto se utiliza el sistema de gestión de base de datos MariaDB (\textit{Ver Figura \ref{diagrama_tablas}}) para gestionar toda la información relacionada con usuarios, personas, casos de desaparición, alertas y grupos de rescate. En ella se han creado las tablas necesarias para le funcionamiento del sistema: `users', `people', `missing people', `reporters', `cases', `found cases', `alerts', `searches', `search participants', `rescue groups' y `group members'. En las figuras \ref{tabla_users}--\ref{tabla_group_members} se muestran capturas de cada una de estas tablas creadas.

    La comunicación entre la aplicación y la base de datos se ha realizado mediante consultas SQL ejecutadas desde el servidor (Node.js + Express). Las operaciones principales implementadas siguen el patrón CRUD (Create, Read, Update, Delete) y se exponen mediante endpoints HTTP: \textbf{GET} para lectura, \textbf{POST} para creación, \textbf{PUT/PATCH} para actualización y \textbf{DELETE} para eliminación.

    A continuación se describen las consultas implementadas para cada tabla:

    \subsubsection{Tabla Users}
    Esta tabla registra la infomración de los usuarios que están asociados a una persona.
    \begin{itemize}
            \item GET que devuelve lista de usuarios con un límite opcional.
            \item GET que obtiene el usuario asociado a un person id.
            \item GET que selecciona un usuario en base a su user id. 
            \item POST que crea un nuevo usuario y devuelve la fila insertada. 
            \item PUT que permite actualizar los campos del usuario indicado, excepto los campos protegidos (user id, person id y created at).   
    \end{itemize}
    \figura{0.7}{img/Resultados/alertresDB/tablaBBDD_Users.jpeg}{Tabla donde se almacena información de los usuarios registrados en la aplicación. \textit{Fuente propia.}}{tabla_users}{keepaspectratio}

    \subsubsection{Tabla People}
        Esta tabla contiene los datos básicos de cualquier persona.
        \begin{itemize}
            \item GET que devuelve la lista de personas con límite opcional.
            \item GET que selecciona una persona en base a su person id.
            \item GET que selecciona a una persona en base a su dni.
            \item POST para crear una nueva persona y devolvuelve la fila insertada.
        \end{itemize}
        \figura{0.7}{img/Resultados/alertresDB/tablaBBDD_People.jpeg}{Tabla donde se almacena información básica de una persona. \textit{Fuente propia.}}{tabla_people}{keepaspectratio}
        
    \subsubsection{Tabla Missing People}
        Que almacena datos más específicos de las personas desaparecidas.
        \begin{itemize}
            \item GET que devuelve la lista de personas desaparecidas con límite opcional.
            \item GET que selecciona una persona desaparecida en base a su missing id.
            \item POST para crear a una nueva persona desaparecida y devuelve la fila insertada.
        \end{itemize}
        \figura{0.7}{img/Resultados/alertresDB/tablaBBDD_MissingPeople.jpeg}{Tabla donde se almacena información específica de las personas desaparecidas. \textit{Fuente propia.}}{tabla_missing_people}{keepaspectratio}
        
    \subsubsection{Tabla Reporters}
        Esta tabla registra a los denunciantes de cada caso de desaparición.
        \begin{itemize}
            \item GET que devuelve la lista de denunciantes con límite opcional.
            \item POST para crear un nuevo denunciante y devolvuelve la fila insertada.
        \end{itemize}
        \figura{0.7}{img/Resultados/alertresDB/tablaBBDD_Reporters.jpeg}{Tabla donde se almacena información de las personas denunciantes. \textit{Fuente propia.}}{tabla_reporters}{keepaspectratio}

    \subsubsection{Tabla Cases}
        Esta tabla recoge la información principal de cada caso de desaparición.
        \begin{itemize}
            \item GET que devuelve la lista de casos con límite opcional.
            \item GET que devuelve la lista de casos con alertas públicas con límite opcional.
            \item GET que obteniene todos los casos de un grupo de rescate.
            \item GET que devuelve los casos en base al estado del caso.
            \item GET que selecciona un caso por su case id.
            \item POST para crear un nuevo caso y devuelve la fila insertada.
            \item PUT que permite actualizar un caso completo excepto los campos protegidos.
        \end{itemize}
        \figura{0.7}{img/Resultados/alertresDB/tablaBBDD_Cases.jpeg}{Tabla donde se almacena información  básica del caso. \textit{Fuente propia.}}{tabla_cases}{keepaspectratio}

    \subsubsection{Tabla Found Cases}
        Esta tabla registra la información post-intervención de un incidente.
        \begin{itemize}
            \item GET que devuelve la lista de casos finalizados con límite opcional.
            \item POST para crear un nuevo caso cerrado y devuelve la fila insertada.
        \end{itemize}
        \figura{0.7}{img/Resultados/alertresDB/tablaBBDD_FoundCases.jpeg}{Tabla donde se almacena la información post-intervención de un caso. \textit{Fuente propia.}}{tabla_found_cases}{keepaspectratio}

    \subsubsection{Tabla Alerts}
        Esta tabla recoge las alertas que se han activado para casos de desaparición.
        \begin{itemize}
            \item GET que devuelve la lista de alertas con límite opcional.
            \item GET que devuelve las alertas públicas y aquellas alertas privadas de un grupo de rescate.
            \item POST para crear una nueva alerta y devuelve la fila insertada.  
        \end{itemize}
        \figura{0.7}{img/Resultados/alertresDB/tablaBBDD_Alerts.jpeg}{Tabla de registro de las alertas activadas. \textit{Fuente propia.}}{tabla_alerts}{keepaspectratio}
        
    \subsubsection{Tabla Searches}
        Esta tabla registra la información de las búsquedas sobre el terreno organizadas para encontrar a una persona desaparecida.
        \begin{itemize}
            \item GET que devuelve la lista de búsquedas con límite opcional.
            \item GET que devuelve las búsquedas que son públicas y aquellas búsquedas privadas de un grupo de rescate.
            \item POST para crear una nueva búsqueda y devuelve la fila insertada.
            \item PUT que permite actualizar los campos de una búsqueda excepto aquellos campos que son protegidos.
        \end{itemize}
        \figura{0.7}{img/Resultados/alertresDB/tablaBBDD_Searches.jpeg}{Tabla de registro de las búsquedas organizadas entorno a un caso. \textit{Fuente propia.}}{tabla_searches}{keepaspectratio}
  
    \subsubsection{Tabla Search Participants}
        Esta tabla agrupa la información de los participantes para cada búsqueda.
        \begin{itemize}
            \item GET que devuelve la lista de participantes con límite opcional.
            \item GET que comprueba si un usuario ya participa en una búsqueda en base al search id y el person id.
            \item POST para crear a un nuevo participante y devuelve la fila insertada.
        \end{itemize}
        \figura{0.7}{img/Resultados/alertresDB/tablaBBDD_SearchParticipants.jpeg}{Tabla donde se almacena información de las personas participantes en una búsqueda específica. \textit{Fuente propia.}}{tabla_search_participants}{keepaspectratio}

    \subsubsection{Tabla Rescue Groups}
        Esta tabla recopila la información de los grupos de rescate.
        \begin{itemize}
            \item GET que devuelve la lista de grupos de rescate con límite opcional.
            \item GET que obtiene el grupo al que pertenece una persona en base a person id. 
            \item GET que selecciona un grupo en base a su group id.
            \item POST para crear un nuevo grupo de rescate y devuelve la fila insertada.
        \end{itemize}
        \figura{0.7}{img/Resultados/alertresDB/tablaBBDD_RescueGroups.jpeg}{Tabla donde se almacena información de los grupos especializados. \textit{Fuente propia.}}{tabla_rescue_groups}{keepaspectratio}

    \subsubsection{Tabla Group Members}
        Esta tabla registra la información de los miembros de cada grupo de rescate.
        \begin{itemize}
            \item GET que devuelve la lista de miembros de grupos de rescate con límite opcional.
            \item GET que obtiene los miembros de un grupo en base a su grupo id.
            \item GET que obtiene la información de un miembro en base a su person id.
            \item POST para crear a un nuevo miembro de un grupo y devuelve la fila insertada.
        \end{itemize}
        \figura{0.7}{img/Resultados/alertresDB/tablaBBDD_GroupMembers.jpeg}{Tabla donde se almacena información de los miembros de un grupo. \textit{Fuente propia.}}{tabla_group_members}{keepaspectratio}
